\chapter{Giới thiệu}

\section{Mục đích tài liệu}

Tài liệu Đặc tả Yêu cầu Phần mềm (Software Requirements Specification -- SRS)
được xây dựng nhằm mô tả đầy đủ các yêu cầu chức năng và phi chức năng của hệ
thống AITasker.

Tài liệu đóng vai trò là cơ sở thống nhất giữa nhóm phát triển, giảng viên hướng
dẫn và các bên liên quan trong suốt quá trình phân tích, thiết kế, triển khai và
kiểm thử hệ thống.

Ngoài ra, tài liệu còn được sử dụng làm cơ sở để:

\begin{itemize}

\item Phân tích và xác định phạm vi hệ thống.

\item Thiết kế kiến trúc phần mềm.

\item Thiết kế cơ sở dữ liệu.

\item Xây dựng giao diện người dùng.

\item Triển khai các chức năng của hệ thống.

\item Kiểm thử và đánh giá chất lượng phần mềm.

\item Bảo trì và mở rộng hệ thống trong tương lai.

\end{itemize}

%=========================================================

\section{Phạm vi hệ thống}

AITasker (AI Marketplace Platform for AI Automation Services) là nền tảng trực
tuyến hỗ trợ kết nối giữa doanh nghiệp có nhu cầu triển khai các giải pháp Trí tuệ
nhân tạo (Artificial Intelligence -- AI) với các chuyên gia AI có đủ năng lực thực
hiện dự án.

Hệ thống hướng đến việc số hóa toàn bộ quy trình thuê chuyên gia AI, từ việc tạo
dự án, tìm kiếm chuyên gia, gửi đề xuất, ký kết hợp đồng, quản lý tiến độ, thanh
toán cho đến đánh giá sau khi hoàn thành dự án.

Ngoài chức năng Marketplace truyền thống, hệ thống còn tích hợp các công cụ AI
nhằm hỗ trợ người dùng trong quá trình xây dựng yêu cầu dự án, tư vấn giải pháp,
gợi ý chuyên gia phù hợp và nâng cao trải nghiệm sử dụng.

Đối tượng sử dụng chính của hệ thống bao gồm:

\begin{itemize}

\item Doanh nghiệp.

\item Chuyên gia AI.

\item Quản trị viên hệ thống.

\end{itemize}

%=========================================================

\section{Đối tượng sử dụng tài liệu}

Tài liệu này được xây dựng dành cho các đối tượng sau:

\begin{itemize}

\item Giảng viên hướng dẫn và hội đồng đánh giá.

\item Nhóm phát triển phần mềm.

\item Kiểm thử viên.

\item Nhà phân tích hệ thống.

\item Người bảo trì hệ thống.

\item Các bên liên quan đến dự án.

\end{itemize}

%=========================================================

\section{Mục tiêu của hệ thống}

Hệ thống AITasker được xây dựng nhằm đạt được các mục tiêu sau:

\begin{enumerate}

\item Kết nối doanh nghiệp với các chuyên gia AI trên cùng một nền tảng.

\item Chuẩn hóa quy trình thuê chuyên gia AI.

\item Quản lý xuyên suốt vòng đời của dự án AI.

\item Hỗ trợ doanh nghiệp xây dựng yêu cầu dự án bằng AI.

\item Tăng tính minh bạch trong quá trình hợp tác giữa các bên.

\item Hỗ trợ thanh toán an toàn thông qua cơ chế quản lý hợp đồng và giao dịch.

\item Cung cấp hệ thống đánh giá nhằm nâng cao chất lượng dịch vụ.

\item Hỗ trợ quản trị viên quản lý toàn bộ hoạt động của hệ thống.

\end{enumerate}

%=========================================================

\section{Định nghĩa, thuật ngữ và từ viết tắt}

\begin{table}[H]
\centering
\caption{Các thuật ngữ và từ viết tắt}
\begin{tabular}{|p{4cm}|p{10cm}|}
\hline
\textbf{Thuật ngữ} & \textbf{Ý nghĩa} \\ \hline

AI &
Artificial Intelligence (Trí tuệ nhân tạo). \\ \hline

Marketplace &
Nền tảng kết nối giữa doanh nghiệp và chuyên gia AI. \\ \hline

Enterprise &
Doanh nghiệp có nhu cầu thuê chuyên gia AI. \\ \hline

Expert &
Chuyên gia AI cung cấp dịch vụ cho doanh nghiệp. \\ \hline

Proposal &
Đề xuất của chuyên gia gửi cho dự án. \\ \hline

Project &
Dự án AI do doanh nghiệp tạo. \\ \hline

Contract &
Hợp đồng giữa doanh nghiệp và chuyên gia. \\ \hline

Milestone &
Mốc công việc trong hợp đồng. \\ \hline

Payment &
Khoản thanh toán theo từng mốc công việc. \\ \hline

Review &
Đánh giá giữa doanh nghiệp và chuyên gia sau khi hoàn thành dự án. \\ \hline

JWT &
JSON Web Token dùng để xác thực người dùng. \\ \hline

REST API &
Giao diện lập trình ứng dụng theo kiến trúc REST. \\ \hline

ORM &
Object Relational Mapping. \\ \hline

SRS &
Software Requirements Specification. \\ \hline

\end{tabular}
\end{table}

%=========================================================

\section{Tài liệu tham khảo}

Trong quá trình phân tích và thiết kế hệ thống, nhóm tham khảo các tài liệu sau:

\begin{enumerate}

\item IEEE Std 830-1998 -- Software Requirements Specification.

\item IEEE Std 29148-2018 -- Requirements Engineering.

\item UML 2.5 Specification.

\item Tài liệu FastAPI.

\item Tài liệu SQLAlchemy.

\item Tài liệu PostgreSQL.

\item Tài liệu ReactJS.

\item Tài liệu JWT (JSON Web Token).

\item Google Gemini API Documentation.

\end{enumerate}

%=========================================================

\section{Cấu trúc tài liệu}

Tài liệu SRS được tổ chức thành các chương như sau:

\begin{itemize}

\item Chương 1: Giới thiệu.

\item Chương 2: Mô tả tổng quan hệ thống.

\item Chương 3: Đặc tả yêu cầu chức năng.

\item Chương 4: Yêu cầu phi chức năng.

\item Chương 5: Thiết kế hệ thống.

\item Chương 6: Thiết kế cơ sở dữ liệu.

\item Chương 7: Phụ lục.

\end{itemize}